\chapter*{\centering ЗАКЛЮЧЕНИЕ}
\addcontentsline{toc}{chapter}{ЗАКЛЮЧЕНИЕ}


Таким образом, в процессе выполнения работы были решены следующие
задачи:
\begin{itemizePaper}
    \item Была изучена и реализована модель обучения с 
    дифференциальным слоем. В ходе численного эксперимента 
    подтверждено, что данная архитектура позволяет эффективно 
    прогнозировать температуру сольвуса с высокой точностью, 
    что решает первую поставленную задачу.
    
    \item Для прогнозирования предела прочности была обучена 
    байесовская регуляризованная нейронная сеть (BRANN). 
    Вторая задача решена в не полном объеме.
    
    \item Проведенный сравнительный анализ показал, что 
    воспроизведенные модели демонстрируют результаты, 
    сопоставимые с оригинальными данными из научных статей, 
    что подтверждает их адекватность и применимость в задачах 
    материаловедения.

\end{itemizePaper}

В пути дальнейшей работы входит донастройка параметров сети 
и/или проба другой арихитектуры нейронно сетию

Таким образом, поставленная цель работы достигнута. Результаты
исследования демонстрируют высокий потенциал методов искусственного 
интеллекта для ускорения разработки и оптимизации состава 
жаропрочных никелевых суперсплавов.

